%% A282_Die_Rechenkugel_De.tex
%% FFGFT A-Serie -- A282: Die Rechenkugel
%% Autor: Johann Pascher  --  ORCID 0009-0000-6518-4064
%% Sprache: Deutsch  --  Engine: LuaLaTeX
%% Prüfskript: a282_rechenkugel.py (Checks 1--11)

\documentclass[11pt,a4paper]{article}

% ---- Engine: LuaLaTeX ----
\usepackage{fontspec}
\usepackage{mathtools}
\usepackage{unicode-math}
\usepackage{polyglossia}
\setmainlanguage{german}

\usepackage{microtype}
\emergencystretch=3.5em
\tolerance=2500
\hbadness=3000
\usepackage{enumitem}
\usepackage{booktabs}
\usepackage{longtable}
\usepackage{array}
\usepackage[a4paper,margin=2.6cm,headsep=1.1em,headheight=22pt]{geometry}
\usepackage{fancyhdr}
\usepackage{titlesec}
\usepackage{xcolor}
\usepackage{needspace}
\usepackage{tocloft}

% ---- Fonts ----
\setmainfont{Inter}[
  Scale=1.02,
  Ligatures=TeX ]
\setsansfont{Inter}[Scale=MatchLowercase,Ligatures=TeX]
\setmonofont{JetBrains Mono}[Scale=0.88]
\setmathfont{Libertinus Math}

% ---- Colours ----
\definecolor{accent}{HTML}{1F3A5F}
\definecolor{accent2}{HTML}{6A4E2E}
\definecolor{rulecol}{HTML}{B9C2CE}
\definecolor{linknav}{HTML}{2C5F8A}
\definecolor{linkcite}{HTML}{7A5C2E}

% ---- Schichtmarkierungen (A-Serie) ----
\definecolor{laySet}{HTML}{7A2E2E}
\definecolor{layK}{HTML}{1D5C3A}
\definecolor{layB}{HTML}{1F3A5F}
\definecolor{layS}{HTML}{6A4E2E}
\definecolor{layH}{HTML}{5A2E6A}

\newcommand{\LS}{\textcolor{laySet}{\textbf{\footnotesize\textsf{[SETZUNG]}}}\,}
\newcommand{\LK}{\textcolor{layK}{\textbf{\footnotesize\textsf{[K]}}}\,}
\newcommand{\LB}{\textcolor{layB}{\textbf{\footnotesize\textsf{[B]}}}\,}
\newcommand{\LSk}{\textcolor{layS}{\textbf{\footnotesize\textsf{[S]}}}\,}
\newcommand{\LH}{\textcolor{layH}{\textbf{\footnotesize\textsf{[H]}}}\,}

\newcommand{\pruef}[1]{{\small\itshape(Prüfskript, #1)}}

% ---- Hyperref ----
\usepackage[
  colorlinks=true,
  linkcolor=linknav,
  citecolor=linkcite,
  urlcolor=linknav,
  bookmarksnumbered=true,
  bookmarksopen=true,
  pdfstartview=FitH ]{hyperref}
\hypersetup{
  pdftitle={Die Rechenkugel (A282)},
  pdfauthor={Johann Pascher},
  pdfsubject={Token-Rechnen als Grenzfall der Landauer-Buchhaltung},
  pdfkeywords={Landauer, Abakus, Token, Phasenraum, Wärmebad, FFGFT, A282}
}

% ---- Section styling ----
\titleformat{\section}
  {\color{accent}\normalfont\Large\bfseries}
  {\color{accent2}\thesection}{0.7em}{}
\titleformat{\subsection}
  {\color{accent}\normalfont\large\bfseries}
  {\color{accent2}\thesubsection}{0.6em}{}
\titlespacing*{\section}{0pt}{1.6\baselineskip}{0.7\baselineskip}
\titlespacing*{\subsection}{0pt}{1.1\baselineskip}{0.4\baselineskip}

% ---- Headers ----
\pagestyle{fancy}
\fancyhf{}
\renewcommand{\headrulewidth}{0.4pt}
\renewcommand{\headrule}{\hbox to\headwidth{\color{rulecol}\leaders\hrule height \headrulewidth\hfill}}
\fancyhead[L]{\footnotesize\color{accent}\itshape FFGFT $\cdot$ A282 $\cdot$ Die Rechenkugel}
\fancyhead[R]{\footnotesize\color{accent2}\thepage}

% ---- TOC look ----
\renewcommand{\cftsecfont}{\normalfont}
\renewcommand{\cftsecpagefont}{\normalfont}
\renewcommand{\cftsecleader}{\cftdotfill{\cftdotsep}}
\setlength{\cftbeforesecskip}{3pt}

\setlength{\parskip}{0.35em}
\setlength{\parindent}{0pt}
\linespread{1.06}

\begin{document}

% ===================== TITEL =====================
\begin{titlepage}
\centering
\vspace*{2.2cm}
{\color{rulecol}\rule{0.5\textwidth}{0.6pt}}\\[1.4cm]
{\color{accent}\Huge\bfseries Die Rechenkugel\par}
\vspace{1.0cm}
{\large\itshape Token-Rechnen als Grenzfall\\ der Landauer-Buchhaltung\par}
\vspace{0.6cm}
{\color{accent2}\rule{0.28\textwidth}{0.4pt}}\\[1.0cm]
{\large FFGFT A-Serie $\cdot$ A282\par}
\vspace{0.4cm}
{\normalsize Drittes Stück zu A280 (Buchhaltung) und A281 (Textkritik)\par}
\vfill
{\large Johann Pascher\par}
\vspace{0.3cm}
{\footnotesize ORCID: 0009-0000-6518-4064\par}
\vspace{0.2cm}
{\footnotesize Juli 2026\par}
\vspace{1.4cm}
{\color{rulecol}\rule{0.5\textwidth}{0.6pt}}
\vspace*{1.5cm}
\end{titlepage}

% ===================== TOC =====================
\tableofcontents
\thispagestyle{fancy}
\clearpage

% ===================== KURZFASSUNG =====================
\begingroup\small\noindent
\textbf{Kurzfassung.} Rechnen mit Marken --- Abakuskugeln, Muscheln, Münzen ---
ist der Grenzfall, an dem sich Landauers Argument in seine zwei Hälften zerlegen
lässt. Die \emph{Buchhaltung} gilt dort exakt und ist zum ersten Mal sichtbar: Der
Vorrat, aus dem die Marken kommen und in den sie zurückfließen, ist ein Bad, auf
das man zeigen kann. Die \emph{thermische Umrechnung} $\Delta S \to Q = T\Delta S$
gilt dort nicht: Eine Abakuskugel kostet beim Verschieben rund $5\cdot10^{15}$ mal
die Landauer-Grenze, ihr Positions-Bit trägt $4\cdot10^{-23}$ ihrer eigenen
thermischen Entropie, und ihre Positionsbarriere liegt bei $3{,}6\cdot10^{15}\,kT$.
Damit fällt der Token-Rechner unter genau den Vorbehalt, den Landauer selbst
formuliert und die Rezeption fallen gelassen hat (A281): Die gezählten Zustände
sind kein thermisches Ensemble. Daraus folgt eine Ordnung, die der Intuition
zuwiderläuft: \textbf{Sichtbarkeit der Struktur und Bindekraft der Zahl laufen
gegenläufig.} Je primitiver der Träger, desto klarer die Buchhaltung und desto
irrelevanter die Schranke. Der Grenzübergang zurück ist stetig --- schrumpft die
Marke, bis ihre Barriere gegen $kT$ geht, wird sie zur Brownschen Rechenkugel, und
das ist Béruts Kolloid.
\par\endgroup

\bigskip

\begingroup\small\noindent
\textbf{Schichtmarkierungen.} \LS deklarierter Ausgangspunkt ---
\LB algebraisch/mathematisch aus deklarierten Grundlagen abgeleitet ---
\LK empirisch geprüft bzw.\ aus Primärquellen belegt ---
\LSk Plausibilitätsskizze ---
\LH offene Forschungsfrage.
\par\endgroup

\bigskip

\begingroup\small\noindent
\textbf{Prüfskript.} \texttt{a282\_rechenkugel.py} (Checks 1--11). Sämtliche
Zahlen dieses Dokuments sind dort mit Assertions hinterlegt; die
Parameterwahlen sind im Skript als Setzungen ausgewiesen und frei änderbar.
\par\endgroup

\bigskip

% ===================== 1 =====================
\section{Der Vorrat als sichtbares Bad}

A280 arbeitet mit dem Weltmeer: Ein Glas wird eingeschüttet, das Volumen
verschwindet nicht, aber niemand holt es zurück. Das Bild ist richtig und hat
einen Nachteil --- man muss das Bad erschließen. Niemand sieht die $10^{23}$
Freiheitsgrade, in die das Bit-Volumen kippt.

\LS Beim Rechnen mit Marken ist das anders. Ein Abakus, ein Muschelhaufen, eine
Kasse mit Münzen: Neben dem Arbeitsbereich liegt ein \textbf{Vorrat}. Aus ihm
kommen die Marken, in ihn fließen sie zurück. Der Vorrat ist ein Bad, auf das man
zeigen kann.

\begin{center}
\begin{minipage}{0.80\textwidth}
\ttfamily\small
\begin{tabbing}
\ Arbeitsbereich\ \ \ \ \ \ \ \ \ \ \=\ Vorrat (Bad)\\[0.3em]
+-----------------+\>+-------------------------+\\
|\ o\ o\ .\ o\ .\ .\ o\ .\ |\>|\ o o o o o o o o o o o o |\\
|\ .\ o\ o\ .\ o\ o\ .\ o\ |\ \ ---\textgreater\ \>|\ o o o o o o o o o o o o |\\
+-----------------+\>+-------------------------+\\
\ W = 2\^{}N Konfig.\>\ Marken erhalten, Lage gemischt
\end{tabbing}
\end{minipage}
\end{center}

\LB Die Struktur ist Zug um Zug dieselbe wie in A280:

\begin{itemize}[leftmargin=1.4em,itemsep=0.3em]
  \item \textbf{Gebiet.} Eine Marke auf Position links und dieselbe Marke auf
        Position rechts sind zwei unterscheidbare Makro-Gebiete.
  \item \textbf{Gebietsreduktion.} Das Zurücksetzen des Rechenbretts bildet alle
        Konfigurationen auf eine ab --- alle Marken im Vorrat. Nicht-bijektiv,
        genau wie $\{A,B\}\to Z$.
  \item \textbf{Erhaltung.} Keine Marke verschwindet. Sie verteilen sich auf eine
        Lage, aus der die vorherige Konfiguration nicht mehr ablesbar ist. Das ist
        das Liouville-Analogon --- und es ist hier buchstäblich abzählbar.
  \item \textbf{Inhaltsneutralität.} Der Vorrat weiß nicht, ob eine Marke einen
        Schuldposten, eine Fünf oder gar nichts bedeutet hat. Er nimmt sie gleich
        auf.
\end{itemize}

\LB Bargeld mit Münzen und einer Kasse ist dieselbe Figur, und die Doppelbedeutung
des Wortes „Buchhaltung“ ist hier keine Metapher mehr: Kassenbestand und
Phasenraum-Buchhaltung haben dieselbe Form. Der Tauschhandel ist der Fall ohne
Marken --- dort gibt es keinen Vorrat und entsprechend keine Gebietsreduktion,
sondern nur bijektiven Austausch.

% ===================== 2 =====================
\section{Die Buchhaltung gilt exakt}

\LB Für $N$ Marken mit je zwei Positionen gilt vor dem Rücksetzen
$W_{\text{vorher}} = 2^N$ und danach $W_{\text{nachher}} = 1$, also
\begin{equation}
  \frac{\Delta S}{k_B} \;=\; \ln\frac{W_{\text{vorher}}}{W_{\text{nachher}}}
  \;=\; N \ln 2 .
  \label{eq:kombinatorik}
\end{equation}

Das ist exakt und trägerunabhängig \pruef{Check 10}. Es gilt für Kugeln, für
Muscheln, für Transistoren, für Kolloide. Nichts an
Gleichung~\eqref{eq:kombinatorik} verlangt, dass die Marke klein, leicht oder
thermisch bewegt sei. Die Zählung ist Kombinatorik über Gebiete --- genau das,
was A280 als das Primäre ausweist.

\LB Bis hierher ist der Token-Rechner ein vollständiges Modell der
Landauer-Buchhaltung, und zwar ein besseres als jedes andere, weil man jeden
Schritt sehen kann.

% ===================== 3 =====================
\section{Die Schranke gilt nicht}

\LS Für die folgenden Zahlen werden gesetzt: Abakuskugel der Masse
$m = 0{,}5$\,g, Verschiebeweg $d = 1$\,cm, Reibungszahl $\mu = 0{,}3$,
Umgebungstemperatur $T = 300$\,K, molare Masse $100$\,g/mol und molare Entropie
$50$\,J/(mol$\cdot$K) für einen Festkörper. Alle Setzungen sind im Prüfskript
markiert und änderbar; die Größenordnungen hängen nicht an ihrer genauen Wahl.

\LB \textbf{Die Arbeit je Kugelverschiebung.} Gleitreibung liefert
$W = \mu\,m\,g\,d = 1{,}47\cdot10^{-5}$\,J \pruef{Check 2}. Die Landauer-Grenze
bei 300\,K beträgt $k_B T \ln 2 = 2{,}87\cdot10^{-21}$\,J \pruef{Check 1}. Das
Verhältnis ist
\begin{equation}
  \frac{W_{\text{Kugel}}}{k_B T \ln 2} \;\approx\; 5{,}1\cdot10^{15}
  \qquad \text{\pruef{Check 3}} .
\end{equation}

\subsection{Das Skalenband}

\begin{center}
\small
\begin{tabular}{@{}lll@{}}
\toprule
\textbf{Träger} & \textbf{Kosten je Schaltvorgang} & \textbf{Vielfaches von $k_BT\ln 2$} \\
\midrule
Kolloid, quasistatisch [B4][B5] & ${\sim}3\cdot10^{-21}$\,J & 1--10 \\
CMOS-Transistor (A280) & $10^{-17}$--$10^{-16}$\,J & $3{,}5\cdot10^{3}$ -- $3{,}5\cdot10^{4}$ \\
Abakuskugel & $1{,}5\cdot10^{-5}$\,J & $5{,}1\cdot10^{15}$ \\
\bottomrule
\end{tabular}
\end{center}

\LB Zwischen Kugel und CMOS liegen rund \textbf{elf Dekaden} \pruef{Check 5}. Die
Ordnung ist damit die umgekehrte der naheliegenden Vermutung: Landauer bindet
nicht am primitiven Ende, sondern am kleinen. Je gröber der Träger, desto weiter
liegt die Schranke unter dem tatsächlichen Aufwand.

\subsection{Warum die Kugel nichts zahlt}

Zwei Zahlen erklären den Befund, und beide betreffen die Kugel selbst, nicht die
Buchhaltung.

\LB \textbf{Erstens: Das Positions-Bit ist verschwindend gegen die Eigenentropie
der Kugel.} Eine Kugel von 0,5\,g trägt bei den obigen Setzungen eine Entropie von
$0{,}25$\,J/K, also $1{,}8\cdot10^{22}\,k_B$ \pruef{Check 6}. Das Positions-Bit
trägt $\ln 2 = 0{,}69\,k_B$. Sein Anteil ist
\begin{equation}
  \frac{\ln 2}{S_{\text{Kugel}}/k_B} \;\approx\; 3{,}8\cdot10^{-23}
  \qquad \text{\pruef{Check 7}} .
\end{equation}
Beim Verschieben ändert sich am thermisch zugänglichen Phasenraum praktisch
nichts. Die Gebietsreduktion findet in einer Koordinate statt, die von den
$10^{22}$ thermischen Freiheitsgraden der Kugel abgekoppelt ist.

\LB \textbf{Zweitens: Die Positionsbarriere ist thermisch unüberwindlich.} Sie
beträgt $3{,}6\cdot10^{15}\,k_BT$ \pruef{Check 8}. Eine Arrhenius-Lebensdauer
$\propto \exp(3{,}6\cdot10^{15})$ ist numerisch nicht darstellbar. Es gibt keine
thermische Umbesetzung zwischen „links“ und „rechts“, also auch kein
Gleichgewichtsensemble über diese beiden Zustände.

\LB Damit ist die Kugel der Extremfall der Zeile „emergente Zweiwertigkeit“ aus
A280: Die Zwei ist eine Einteilung, die einem Quasikontinuum von $10^{22}$
Mikrozuständen aufgelegt wird, und sie überlebt nicht deshalb, weil sie
fundamental wäre, sondern weil die Barriere absurd hoch ist.

% ===================== 4 =====================
\section{Die zwei Hälften von Landauers Argument}
\label{sec:zweihaelften}

\LB Der Token-Rechner leistet damit etwas, was kein anderes Modell so sauber
leistet: Er trennt die beiden Schritte, aus denen Landauers Argument besteht, und
zeigt, dass sie unabhängig sind.

\begin{center}
\small
\begin{tabular}{@{}p{0.16\textwidth}p{0.36\textwidth}p{0.40\textwidth}@{}}
\toprule
\textbf{Schritt} & \textbf{Aussage} & \textbf{Voraussetzung} \\
\midrule
Kombinatorik & $\Delta S = k_B \ln(W_{\text{v}}/W_{\text{n}})$ &
Gebiete sind unterscheidbar und abzählbar \\
Umrechnung & $Q = T\,\Delta S$ &
die gezählten Zustände sind thermisch besetzt \\
\bottomrule
\end{tabular}
\end{center}

\LB Der erste Schritt gilt für die Kugel exakt. Der zweite gilt für sie
überhaupt nicht. Die übliche Darstellung führt beide in einem Zug aus und macht
die zweite Voraussetzung damit unsichtbar. Erst wenn ein Träger sie maximal
verletzt, tritt sie hervor.

\begin{quote}
\LB Landauers Schranke ist nicht die Aussage „Gebietsreduktion kostet“. Sie ist
die Aussage „Gebietsreduktion \emph{in thermisch besetzten Zuständen} kostet“.
Der Zusatz ist nicht kosmetisch --- er trägt die ganze Umrechnung.
\end{quote}

\LB Das schärft A280 an einer Stelle nach, an der es genauer sein muss. A280
§6 verlangt für klassische Systeme Unterscheidbarkeit \emph{oberhalb} der
thermischen Auflösungsgrenze. Das ist die Bedingung dafür, dass zwei Gebiete als
zwei zählen. Es ist \emph{nicht} die Bedingung dafür, dass ihre Reduktion Wärme
kostet --- dafür braucht es zusätzlich, dass die Barriere \emph{nicht} beliebig
weit oberhalb $k_BT$ liegt. Beide Bedingungen zeigen in entgegengesetzte
Richtungen, und dazwischen liegt das schmale Band, in dem Landauer bindet.

% ===================== 5 =====================
\section{Landauers eigener Vorbehalt}

\LK Der Befund ist keine Neuentdeckung, sondern der Extremfall einer
Einschränkung, die Landauer 1961 selbst formuliert hat: Seine Rücksetz-Operation
wurde auf ein thermisches Gleichgewichts-Ensemble angewandt; für nicht
thermalisierte Information gilt die Rechnung nicht unmittelbar [B1]. A281
dokumentiert die Passage und weist nach, dass die Rezeption sie fallen gelassen
hat.

\LB Beim Abakus ist der Vorbehalt maximal verletzt, und zwar auf zwei Weisen
zugleich:

\begin{enumerate}[leftmargin=1.6em,itemsep=0.3em]
  \item \textbf{Kein thermisches Ensemble.} Die Verteilung über „links“ und
        „rechts“ entsteht nicht durch thermische Besetzung, sondern durch die
        Hand, die die Kugel gesetzt hat.
  \item \textbf{Vollständiges Wissen.} Wer die Kugeln gesetzt hat, kennt die
        Konfiguration. In der Rückkopplungsschranke (A281) ist damit
        $I = \ln 2$ pro Marke und
        \begin{equation}
          Q_{\min} \;=\; k_B T\,(\ln 2 - I) \;=\; 0
          \qquad \text{\pruef{Check 11}} .
        \end{equation}
\end{enumerate}

\LB Thermodynamisch kostet das Rücksetzen eines bekannten Abakus also nichts. Was
es kostet, ist Reibung --- und die ist eine Ingenieursgröße ohne untere
Naturschranke, beliebig verkleinerbar durch bessere Lager.

% ===================== 6 =====================
\section{Der Grenzübergang: die Brownsche Rechenkugel}

\LB Der Übergang zurück in Landauers Geltungsbereich ist stetig und lässt sich
ausrechnen. Setzt man die Positionsbarriere gleich $k_BT$, so folgt für die
kritische Markenmasse
\begin{equation}
  \mu\,m\,g\,d \;=\; k_B T
  \quad\Longrightarrow\quad
  m_{\text{krit}} \;=\; \frac{k_B T}{\mu\,g\,d} \;\approx\; 1{,}4\cdot10^{-19}\,\text{kg}
  \qquad \text{\pruef{Check 9}} .
\end{equation}

\LK Zum Vergleich: Ein Silika-Kolloid von 2\,\textmu m Durchmesser wiegt
$8{,}4\cdot10^{-15}$\,kg, also rund $6\cdot10^{4}$ mal mehr \pruef{Check 9}. Es
hält seine Position aber nicht durch Reibung, sondern durch eine optische Falle,
und deren Barrierenhöhe ist auf wenige $k_BT$ einstellbar. Genau das macht es zur
Brownschen Rechenkugel --- und genau deshalb ist Béruts Aufbau [B4] der einzige,
in dem die Schranke tatsächlich bindet.

\begin{quote}
\LB Béruts Kolloid ist kein Modell des Bits. Es ist eine Abakuskugel, die klein
genug ist, dass das Bad sie schiebt.
\end{quote}

\LB Damit schließt sich die Reihe: Muschel, Münze, Abakuskugel, Transistor,
Kolloid. Es ist eine Skala, kein Sprung. Was sich entlang der Skala ändert, ist
nicht die Buchhaltung --- sie ist überall dieselbe --- sondern das Verhältnis der
Positionsbarriere zu $k_BT$.

% ===================== 7 =====================
\section{Zwei Achsen, gegenläufig}

\LB Daraus folgt die Ordnung, die dem Dokument seinen Ertrag gibt:

\begin{center}
\small
\begin{tabular}{@{}lll@{}}
\toprule
\textbf{Träger} & \textbf{Sichtbarkeit der Struktur} & \textbf{Bindekraft der Schranke} \\
\midrule
Abakus, Muscheln, Münzen & vollständig (Bad zeigbar) & keine ($10^{15}$ darüber) \\
CMOS & gering (Bad ist das Gitter) & schwach ($10^{4}$ darüber) \\
Kolloid in optischer Falle & keine (Bad nicht trennbar) & voll (Faktor 1--10) \\
\bottomrule
\end{tabular}
\end{center}

\begin{quote}
\LB Auf der primitivsten Ebene ist die Buchhaltung maximal sichtbar und die
Schranke maximal irrelevant. Sichtbarkeit der Struktur und Bindekraft der Zahl
sind zwei verschiedene Achsen, und sie laufen gegenläufig.
\end{quote}

\LB Das erklärt zugleich, warum die Anschauung in dieser Sache regelmäßig in die
Irre führt. Wer sich Landauer plausibel machen will, greift zum groben Bild ---
zur Kugel, zur Schublade, zum Glas Wasser --- und landet damit genau dort, wo die
Schranke nichts mehr bedeutet. Das Bild trägt die Struktur und verliert die Zahl.

% ===================== 8 =====================
\section{Der Markenvorrat als eigene Kostenkategorie}

\LB Ein Punkt bleibt, den die thermodynamische Buchhaltung nicht erfasst und der
am Token-Rechner besonders deutlich wird: \textbf{Ohne Marken kann nicht gerechnet
werden.} Der Vorrat ist eine erschöpfbare Ressource. Ein Abakus mit zu wenigen
Kugeln rechnet nicht langsamer, sondern gar nicht.

\LB Das ist eine \emph{Materialschranke}, keine thermodynamische. Sie hat keine
Untergrenze in $k_BT$, sie hat eine Untergrenze in der Stückzahl. Sie ist
verwandt mit dem, was A280 §7.3 aufmacht --- der Rechner zahlt dafür, dass er
existiert, nicht nur dafür, dass er rechnet ---, aber sie ist nicht dasselbe: Dort
geht es um einen laufenden Entropiestrom, hier um einen stehenden Bestand.

\LSk Beim Halbleiter fällt diese Kategorie aus dem Blick, weil der Vorrat in
Silizium eingegossen ist: Die Gebiete werden bei der Fertigung angelegt und nicht
im Betrieb entnommen. Der Bestand ist da, also wird er nicht gezählt. Am Abakus
ist er sichtbar, weil man ihn aufbrauchen kann.

\LH \textbf{Offene Frage.} Ob sich Bestand und Strom in einer gemeinsamen
Buchhaltung führen lassen --- Marken\-vorrat und Entropiestrom als zwei Bilanzen
eines Rechners --- ist offen. Die Frage ist dieselbe wie am Ende von A281, von
der Materialseite gestellt.

% ===================== 9 =====================
\section{Einordnung im Korpus}

\LB Die drei Stücke verhalten sich zueinander so:

\begin{itemize}[leftmargin=1.4em,itemsep=0.3em]
  \item \textbf{A280} führt die Buchhaltung: Gebiet, Gebietsreduktion, Liouville,
        Bad, Schranke. Physikalischer Kern.
  \item \textbf{A281} führt die Textkritik: was Landauer 1961 tatsächlich beweist,
        sein eigener Ensemble-Vorbehalt, die Rezeptionsgeschichte, die Sprache.
  \item \textbf{A282} führt den Grenzfall vor, an dem beide Befunde zusammenlaufen
        und sichtbar werden: Der Token-Rechner erfüllt A280s Buchhaltung exakt und
        verletzt A281s Ensemble-Bedingung maximal.
\end{itemize}

\LB Die Wörterbuch-Zuordnung aus A281 §1.1 gilt unverändert. In der Sprache dieses
Dokuments: Marke $=$ Träger $=$ Gebiet; Vorrat $=$ Bad; Rücksetzen des
Rechenbretts $=$ Träger-Operation $=$ Gebietsreduktion.

\LB \textbf{FFGFT-Berührung: keine.} Wie A280 und A281 lebt auch dieses Dokument
auf der Ebene der statistischen Mechanik; das $\xi$-Schema wird nicht berührt.

% ===================== 10 =====================
\section{Schichtstatus-Zusammenfassung}

\begin{center}
\small
\begin{longtable}{@{}p{0.72\textwidth}l@{}}
\toprule
\textbf{Aussage} & \textbf{Status} \\
\midrule
\endfirsthead
\toprule
\textbf{Aussage} & \textbf{Status} \\
\midrule
\endhead
Der Markenvorrat ist ein zeigbares Bad-Analogon & \LS \\
Parameter der Abakuskugel (Masse, Weg, Reibung, $T$) & \LS \\
Rücksetzen des Rechenbretts ist eine nicht-bijektive Gebietsabbildung & \LB \\
Marken bleiben erhalten (Liouville-Analogon) & \LB \\
$\Delta S/k_B = N\ln 2$, trägerunabhängig exakt & \LB \\
$k_BT\ln 2 = 2{,}87\cdot10^{-21}$\,J bei 300\,K & \LK \\
Kugelverschiebung $1{,}47\cdot10^{-5}$\,J & \LB \\
Kugel $\approx 5{,}1\cdot10^{15}\times$ Landauer & \LB \\
Skalenband Kolloid / CMOS / Kugel & \LK \\
Elf Dekaden zwischen Kugel und CMOS & \LB \\
Eigenentropie der Kugel $1{,}8\cdot10^{22}\,k_B$ & \LB \\
Positions-Bit $= 3{,}8\cdot10^{-23}$ der Eigenentropie & \LB \\
Positionsbarriere $3{,}6\cdot10^{15}\,k_BT$ & \LB \\
Kombinatorik und thermische Umrechnung sind unabhängig & \LB \\
Die Umrechnung setzt thermisch besetzte Zustände voraus & \LB \\
Landauers Ensemble-Vorbehalt (Originalstelle) & \LK \\
Bekannter Abakus: $I=\ln 2$, also $Q_{\min}=0$ & \LB \\
Kritische Markenmasse $1{,}4\cdot10^{-19}$\,kg & \LB \\
Silika-Kolloid 2\,\textmu m: $8{,}4\cdot10^{-15}$\,kg & \LK \\
Béruts Kolloid als Brownsche Rechenkugel & \LB \\
Sichtbarkeit und Bindekraft laufen gegenläufig & \LB \\
Markenvorrat als Materialschranke & \LB \\
Warum die Kategorie beim Halbleiter aus dem Blick fällt & \LSk \\
Gemeinsame Bilanz von Bestand und Strom & \LH \\
\bottomrule
\end{longtable}
\end{center}

\textbf{Bilanz:} 15$\times$ \LB $\cdot$ 4$\times$ \LK $\cdot$ 2$\times$ \LS
$\cdot$ 1$\times$ \LSk $\cdot$ 1$\times$ \LH $=$ 23 Einträge.

% ===================== BELEGE =====================
\clearpage
\phantomsection
\addcontentsline{toc}{section}{Belege}
\section*{Belege}

\begingroup
\small
\begin{description}[leftmargin=3.2em,style=sameline,itemsep=0.3em]

\item[{[B1]}] R. Landauer, \emph{Irreversibility and Heat Generation in the
Computing Process}, IBM J.\ Res.\ Dev.\ \textbf{5}, 183 (1961).

\item[{[B2]}] Liouville-Theorem: Standardresultat der Hamiltonschen Mechanik;
z.\,B.\ Landau/Lifschitz, \emph{Mechanik}, §46.

\item[{[B3]}] C.~H. Bennett, \emph{The Thermodynamics of Computation --- a
Review}, Int.\ J.\ Theor.\ Phys.\ \textbf{21}, 905 (1982).

\item[{[B4]}] A. Bérut et al., \emph{Experimental verification of Landauer's
principle}, Nature \textbf{483}, 187 (2012).

\item[{[B5]}] Y. Jun, M. Gavrilov, J. Bechhoefer, \emph{High-Precision Test of
Landauer's Principle in a Feedback Trap}, Phys.\ Rev.\ Lett.\ \textbf{113},
190601 (2014).

\item[{[B6]}] T. Sagawa und M. Ueda, \emph{Nonequilibrium Thermodynamics of
Feedback Control}, Phys.\ Rev.\ E \textbf{85}, 021104 (2012).

\item[{[B7]}] Semiklassische Zustandszählung $\Gamma/h^f$ und Entropie von
Festkörpern: Standard; z.\,B.\ Landau/Lifschitz, \emph{Statistische Physik}, §7
und §64.

\item[{[B8]}] CODATA-Werte 2018; $k_B = 1{,}380649\cdot10^{-23}$\,J/K seit der
SI-Revision 2019 exakt definiert.

\end{description}
\endgroup

\vspace{1em}
\noindent
{\small\itshape Die Zahlenwerte dieses Dokuments stammen aus dem Prüfskript
\texttt{a282\_rechenkugel.py}; die dort gesetzten Parameter sind Modellwahlen und
keine Messwerte. Die Größenordnungen sind gegen ihre Variation unempfindlich.}

\end{document}
